\section{Procedimiento}

Se utilizará CRISP-DM porque relaciona las necesidades del negocio con la preparación de datos, el modelado, la evaluación y el despliegue. El procedimiento comprende las siguientes fases:

\begin{enumerate}[label=\textbf{\arabic*.}]
  \item \textbf{Comprensión del negocio.} Precisar las necesidades de rentabilidad, pronóstico y análisis de ciclos; definir indicadores y criterios de aceptación con apoyo de especialistas del sector.
  \item \textbf{Comprensión de los datos.} Descargar y documentar los microdatos INEI y las operaciones Aduanet de la subpartida 0804400000; identificar claves, unidades y cobertura temporal.
  \item \textbf{Preparación.} Eliminar duplicados, validar dominios, tratar ausentes, calcular precio FOB unitario, agregar la serie y construir rezagos, temporada, productividad, costos y margen.
  \item \textbf{Modelado.} Ajustar Random Forest, HistGradientBoosting, ElasticNet, regresión logística y Prophet mediante divisiones reproducibles y parámetros documentados; construir un ensamble FOB con las salidas de los cuatro modelos de regresión y pronóstico.
  \item \textbf{Evaluación.} Comparar MAE, RMSE y MAPE por horizonte; revisar residuos, estabilidad e interpretabilidad; seleccionar el modelo que supere la línea base con comportamiento consistente.
  \item \textbf{Desarrollo.} Implementar servicios en FastAPI y una interfaz React o React Native con acceso, pronóstico, histórico, simulación, recomendación, reportes y notificaciones.
  \item \textbf{Validación.} Validar instrumentos con jueces, ejecutar la prueba piloto, analizar confiabilidad y comparar las mediciones pretest y postest.
  \item \textbf{Despliegue y documentación.} Preparar una versión de prueba y documentar código, diccionario de datos, parámetros, resultados y limitaciones.
\end{enumerate}

\begin{figure}[H]
\caption{Procedimiento de la investigación basado en CRISP-DM}
\centering
\begin{tikzpicture}[node distance=7mm, every node/.style={font=\small}, box/.style={draw, rounded corners, align=center, text width=3.7cm, minimum height=9mm}, arr/.style={-{Latex[length=2mm]}, thick}]
\node[box] (n1) {Comprensión del negocio};
\node[box, below=of n1] (n2) {Comprensión de los datos};
\node[box, below=of n2] (n3) {Preparación};
\node[box, below=of n3] (n4) {Modelado};
\node[box, below=of n4] (n5) {Evaluación};
\node[box, below=of n5] (n6) {Desarrollo y despliegue};
\draw[arr] (n1)--(n2);
\draw[arr] (n2)--(n3);
\draw[arr] (n3)--(n4);
\draw[arr] (n4)--(n5);
\draw[arr] (n5)--(n6);
\draw[arr, bend left=50] (n5.east) to (n2.east);
\end{tikzpicture}
\end{figure}
